\documentclass[10pt,a4paper]{amsart}
\usepackage{amsmath,amssymb,amsthm,graphicx}
\usepackage{amsmath,amssymb,amsthm,enumerate,mathabx}
\newtheorem{Theorem}{Theorem}[section]
\newtheorem{Proposition}[Theorem]{Proposition}
\newtheorem{Lemma}[Theorem]{Lemma}
\newtheorem{Corollary}[Theorem]{Corollary}
\newtheorem{Theorem1}{Theorem}
 
\theoremstyle{definition}
\newtheorem{Definition}[Theorem]{Definition}
\newtheorem{Fact}[Theorem]{Fact}
\newtheorem{Conjecture}[Theorem]{Conjecture}
\newtheorem{Example}[Theorem]{Example}
\newtheorem{Remark}[Theorem]{Remark}
\newtheorem{Question}{Question}
\newtheorem{Claim}{Claim}

\DeclareMathOperator{\cl}{cl_S}
\DeclareMathOperator{\dcl}{dcl}
\DeclareMathOperator{\tp}{tp}
\DeclareMathOperator{\Aut}{Aut}
\DeclareMathOperator{\conv}{conv} 
\DeclareMathOperator{\dom}{dom}
\DeclareMathOperator{\acc}{acc}
\title{Realizing prescribed countable discrete orders as the locus of a simple type}
\author{Predrag Tanović}
\date{arxiv:2212.13605v3; CC BY 4.0}
\begin{document}
\maketitle
\setcounter{Theorem1}{1}
\noindent\textit{Editorial scope.} Count only original Theorem2. All source-specific rigidity, uniqueness, and minima proofs are part of the delivered solution chain; they are not granted as free prerequisites. The author credit to D. Ilic at source216 and original bibliography entry are retained. No source mathematics has been changed. The earlier r001 role classification is preserved separately.\par\medskip
\noindent\textit{Allowed standard prerequisites (editorial).}\begin{itemize}
\item Compactness, countable omitting-types theorem, saturated homogeneous monster models. Extend partial definable functions/orders and omit a nonprincipal type over a countable set of named elements. All special-type rigidity and completion claims are supplied from the author.
\item Continuum many isomorphism types of countable discrete endpoint-free linear orders. Infer the final countable-model cardinality after exact locus realization.
\end{itemize}\medskip
Throughout the paper, $L$ will be an at most countable  language and $T$  an arbitrary (possibly multi-sorted), complete,  first-order theory with infinite models;  $I(\aleph_0,T)$ will denote the number of countable models of $T$, up to 
isomorphism.  We say that $T$  admits an infinite,  discrete, linear order, if such a parametrically definable order $(D,<)$ can be found in some   $\aleph_0$-saturated model $M\models T$ with $D\subseteq M^n$; by adding an extra sort, if necessary, we may always assume $D\subseteq M$. 

Let $(P,<)$ be a linear order and $X,Y\subseteq P$. $X<Y$ means that $x<y$ for all $x\in X$ and $y\in Y$.
$x<Y$ means  $\{x\}<Y$.  By $[x,y]$ we will denote the closed interval $\{z\in P\mid x\leq z\leq y\}$.  $X$ is an {\it initial part} if $x\in X$ and $y<x$ imply $y\in X$; $X$ is {\it convex} if $x,y\in X$ and $x<y$ imply $[x,y]\subseteq X$.
The element $y\in P$ is an immediate successor of $x\in P$, denoted by $S(x)=y$, if $x<y$ and no element of $P$ is strictly between $x$ and $y$. We will view $S$ as a partial function on $P$ and $S^m$ will be the $m$-th iterate of $S$; since $S$ is one-to-one, $S^m$ is a well-defined partial function for all  integers $m$. The minimal element of $X$ (or the left end of $X$), if it exists, is denoted by $\min X$; similarly for $\max X$. $(P,<)$ is a discrete order if the induced order-topology is discrete; equivalently: $S(x)$ is defined for all elements of $P$ except $\max P$ (if it exists) and  $S^{-1}(x)$ is defined for all elements of $P$ except $\min P$ (if it exists). 

Our model theoretic notation is mainly standard. 
We work within a fixed  monster model $M\models T$.  $a,b,c,a'...$ denote elements,  $\bar a,\bar b,...$ tuples of elements, and $A,B,...$ small (usually countable)  subsets of $M$; the letters $D,P,X,...$ may denote  any sets. Definability means without parameters and by partial types we will mean those that are closed under finite conjunctions. If $D\subseteq M$ and $\theta(\bar x)$ is a $L_M$-formula,  then  
 $\theta(D)=\{\bar d\in D^{|\bar x|}\mid M\models \theta(\bar d)\}$ is the set of realizations of  $\theta(\bar x)$ in $D$; similarly for partial  types. $\dcl(A)$ denotes the definable closure of $A$ in $M$: $a\in\dcl(A)$ means that there is a $L_A$-formula whose only realization in $M$ is $a$.  
$S_n(A)$ denotes the space of all complete $n$-types with parameters from $A$ endowed with the formula topology (basic clopen sets are $U_{\phi}=\{q\mid \phi(\bar x)\in q\}$ $\phi(\bar x)\in L_A$); 
$S_n(T):=S_n(\emptyset)$). The theory $T$ is {\it small} if   $\bigcup_{n\in\mathbb N}S_n(T)$ is countable. If $T$ is not small,  then  $I(\aleph_0,T)=2^{\aleph_0}$. 

Let $\Pi(\bar x)$ be a partial type over $A$ and $X\subseteq \Pi(M)$. We say that the set $X$ is relatively $A$-definable within $\Pi(M)$ if there is a $L_A$-formula $\theta(x)$ with $X=\Pi(M)\cap \theta(M)$; in this case we say that $\theta(x)$ relatively defines $X$ within $\Pi(M)$. 
If the sets $X_1,X_2$ are type-definable over $A$, then the function $f: X_1\to X_2$  is relatively $A$-definable if its graph is such.


\begin{Fact}\label{Fact1}
(a) Any relatively definable linear order has a definable extension. More precisely, if $\Pi(x)$ is a partial type over $A$ and $\phi(x,y)$  an  $L_A$-formula which relatively defines a linear order on $\Pi(M)$, then $\phi(x,y)$ defines a linear order on $\theta(M)$ for some  formula $\theta(x)\in \Pi(x)$. 


(b)  Any relatively definable function has a definable extension. 
\end{Fact}
\begin{proof}
Routine compactness, we sketch only proof of part (a). Define:
$$\lambda(x,y,z):=\lnot\phi(x,x)\land (x\neq y\rightarrow \lnot(\phi(x,y)\leftrightarrow\phi(y,x)))\land (\phi(x,y)\land\phi(y,z)\rightarrow\phi(x,z)).$$
Since the formula $\phi(x,y)$ defines a linear order on $\Pi(M)$ we have:
$$\Pi(x)\cup\Pi(y)\cup\Pi(z)\vdash \lambda(x,y,z) .$$
By compactness there exists a formula $\theta(x)\in \Pi(x)$ with  $ \{\theta(x),\theta(y),\theta(z)\}\vdash\lambda(x,y,z)$.  
\end{proof}

Let  $\pi(x)$ be a partial type over $A$  and  $C\subseteq \dcl(A)$. We say that $\pi(x)$ is {\it finitely satisfied} in $C$ if every  formula of $\pi(x)$ is satisfied by some element of $C$. A complete type $p\in S_1(A)$ is finitely satisfied in $C$ if and only if $p$ is in the topological closure of $\mathcal C=\{\tp(c/A)\mid c\in C\}$; here, either  $p\in\mathcal C$, or $p$ is an accumulation point of $\mathcal C$; in the latter case,   we say that $p$ is a $C$-type. If the set $A$ is countable, then the space $S_1(A)$ is compact and separable, so whenever $C$ is infinite then a $C$-type $p\in S_1(A)$ exists. $p$ is the unique $C$-type in $S_1(A)$ if and only if for all $L_A$-formulae $\phi(x)$, exactly one of the sets $\phi(C)$ and $\lnot\phi(C)$ is infinite. 

\begin{Definition}\label{Definition_simple}
A complete type $p\in S_1(A)$ is {\it simple} if there is an infinite set $C\subset\dcl(A)$ and an $A$-definable linear order $(D,<)$ such that the following two conditions hold:

(1) $p$ is the unique $C$-type in $S_1(A)$;

(2) $C$ is an initial part of $D$ ordered in the order-type $\boldsymbol{\omega}$.
\end{Definition}

\begin{Remark}
(i) Condition (1) in the previous definition may be replaced by the following, equivalent one \ : 
  for all $\phi(x)\in p$ the set $\lnot \phi(C)$ is finite. 

\smallskip
(ii)  If $p\in S_1(T)$ is a simple type, as witnessed by $C$ and $(D,<)$, then $p(M)\subset D$ holds: Indeed, any formula defining $D$, say $x\in D$, is satisfied by all elements of $C$, so, since $p$ is a $C$-type, $(x\in D)\in p$ and hence $p(M)\subseteq D$. 

\smallskip (iii) Observe that we did  not require in the definition that the witnessing order $(D,<)$ is discrete.  However, by Lemma \ref{Lemma_basic_simple}(b)  a discrete witness will always exists. 
\end{Remark}


\begin{Theorem1}
Suppose that $p\in S_1(T)$ is a simple type, as witnessed by $(C,<)$ and $(D,<)$. Then for every countable, discrete linear order without end-points $\mathbb L$, there is a countable model $M_{\mathbb L}\models T$ such that $(p(M_{\mathbb L}),<)\cong \mathbb L$. In particular,   $I(\aleph_0,T)=2^{\aleph_0}$. 
\end{Theorem1}

\begin{Lemma}\label{Lemma_basic_simple}
If   $p\in S_1(A)$ is a simple type, as witnessed by $C\subset\dcl(A)$ and $(D,<)$, then:

(a) $P(M):=C\cup p(M)$ is an initial part  of $(D,<)$ that is type-definable over $A$.

(b) There exists an  $A$-definable, discretely ordered, initial part $D_0\subseteq D$  with $P(M)\subseteq D_0$. In particular, $(p(M),<)$ is a discrete order.  

(c) Every nonempty, relatively $M$-definable subset of  $P(M)$ has a  minimum.  

(d) If $C$ and $(D_1,<_1)$ also witness the simplicity of $p$ such that $<_1$ and $<$ agree on $C$, then they agree on $P(M)$, too. 
\end{Lemma}
\begin{proof}To simplify the notation, assume $A=\emptyset$ and  write $x<y$ instead of $x\in D\land y\in D\land x<y$. 

(a) Since $C$ is an initial part of $(D,<)$,    $C<p(M)$ holds. 
To prove that $P(M)$ is an initial part of $D$, it suffices to show that $a\in p(M)$ and $b< a$ together imply $b\in P(M)$. Let $\phi(x)\in \tp(b)$. 
Then  $\exists y(\phi(y)\land y< x)\in\tp(a)=p$. Since $p$ is   a $C$-type, there are $c\in C$ and $b'\in D$ with $\models\phi(b')\land b'< c$.
Then $b'< c$ implies $b'\in C$, so   $\phi(x)$ is realized by $b'\in C$. We have just shown that every formula $\phi(x)\in \tp(b)$ is realized in $C$, so the type $\tp(b)$ is finitely satisfied in $C$; 
If $\tp(b)$ is algebraic  then $b\in C$; otherwise, $\tp (b)$ is a $C$-type, so $\tp(b)=p$ holds by the uniqueness of $p$.
In both cases we have $b\in P(M)$, so $P(M)$ is an initial part of $D$. Then it is is easy to see that $P(M)$ is defined by  $P(x):=\{\exists y(\theta(y) \land   x<y)\mid \theta(y)\in p(y)\}$. 

\smallskip
(b) Let $c_0=\min C$ and let $\phi(x)$ be a formula expressing that the interval $[c_0,x]$ is discretely ordered by $<$.  For all $c\in C$ the interval $[c_0,c]$ is finite, so $\models  \phi(c)$ holds. Since $p$ is a $C$-type, $\phi(x)\in p$. Let $D_0$ be the union of all intervals $[c_0,a]\subset D$ where $a\in\phi(M)$.  

\smallskip
(c) Let $\emptyset\neq D_1\subset P(M)$  where $D_1$ is relatively defined by $\theta(x,\bar b)$ ($\bar b\in M$). Pick $a\in p(M)$ with $[c_0,a]\cap D_1\neq\emptyset$.  Let $\phi(x,\bar y)$ be a formula expressing that the set
$[c_0,x]\cap\theta(M,\bar y)$ is nonempty and has minimum. 
In particular, $\models\phi(a,\bar b)$ holds. For all $c\in C$ the interval $[c_0,c]$ is finite, so $\models \forall \bar y \left(\exists z(c_0\leq z\leq c\land \theta(z,\bar y))\rightarrow \phi(c,\bar y)\right)$. Since $p$ is a $C$-type and $a\in p(M)$, 
\begin{center} 
$\models \forall \bar y \left(\exists z(c_0\leq z\leq a\land \theta(z,\bar y))\rightarrow \phi(a,\bar y)\right)$.  
\end{center}
In particular,   $\models  \exists z(c_0\leq z\leq a\land \theta(z,\bar b))\rightarrow \phi(a,\bar b)$.  Now  $[c_0,a]\cap D_1\neq \emptyset$ implies $\models  \exists z(c_0\leq z\leq a\land \theta(z,\bar b))$, so $\models\phi(a,\bar b)$ and $\min ([c_0,a]\cap D_1)$ exists. Hence $\min D_1$ exists, too.

\smallskip
(d) Suppose not. Then for all $a\in p(M)$ the set $D(a)$ defined by $x<a\land x>_1 a$ is nonempty so, by part (c), $b=\min_{<} D(a)$ is well-defined. By part (a), $C<_1p(M)$ implies $b\notin C$, so  $b<a$ implies $b\in p(M)$. In particular, $D(b)\neq\emptyset$ and for any $b'\in D(b)$ we have $b'<b<a\land b'>_1b>_1 a$, contradicting the minimality of $b$ in $D(a)$. 
\end{proof}

In Section 2 we will work  within the order $(P(M),<)$, which is an initial part of any  order $(D,<)$ that witnesses the simplicity of $p$. By Lemma \ref{Lemma_basic_simple}(d) the order $(P(M),<)$ is determined by   $(C,<)$, so we will simply say that the order $(C,<)$ witnesses the simplicity of the type $p$; $(D,<)$ will be any discrete order, whose initial part is $C$. Note that the  successor function  of $(P(M),<)$ agrees  with the restricted  successor function of $(D,<)$. 

In this section we will simplify arguments from Section 2 of \cite{DI} and use them to prove Theorem 2. The crux of the proof is in the next lemma (Lemma 2.6 in \cite{DI}). 


\begin{Lemma}\label{Lemma_f_is_Sn}
Suppose that  $p\in S_1(T)$ is a simple type, as witnessed by $(C,<)$ and $(D,<)$.  If $f:p(M)\to p(M)$ is a relatively definable function, then $f=S^k$ for some $k\in\mathbb Z$, where $S$ is the successor function on $(p(M),<)$.
\end{Lemma}
\begin{proof}Let $c_0=\min C$, $c_n=S^n(c_0)$ for $n\in\mathbb N$, and  suppose that $(D,<)$ is discrete.   
First, we will prove that $f$ is an automorphism of $(p(M),<)$. Choose a definable function $\hat f:D_1\to D_2$ extending $f$. Since $p(x)\cup\{\hat f(x)=a\}$ is consistent for some (equivalently all) $a\in p(M)$ and $M$ is $\aleph_0$-saturated, $f$ is surjective.
To verify the injectivity of $f$, let us consider the set $\hat f^{-1}(\{f(a)\})\cap [c_0,a]\subset p(M)$. It is an $a$-definable subset of $p(M)$ so, by Lemma \ref{Lemma_basic_simple}(c), has minimum, say $a'$; clearly,   $a'=\min(\hat f^{-1}(\{f(a')\})\cap [c_0,a'])$. Since  $a'\in p(M)$, the formula  $x=\min(\hat f^{-1}(\{f(x)\})\cap [c_0,x])$ belongs to $p(x)$; it follows that $f$ is injective. 
In particular, the inverse function $f^{-1}$ is relatively definable, so after possibly shrinking $D_1$ and $D_2$, we will assume that $\hat f$ is a bijection.
It remains to show that $f$ (equivalently $f^{-1}$) is order-preserving. 
Note that at least one of the formulae $\hat f(x)\leq x$ and $\hat f^{-1}(x)\leq x$ belongs to $p(x)$. We will from now on assume $(\hat f(x)\leq x)\in p(x)$, the proof in the other case is similar.
Suppose, by way of contradiction, that $f$ is not order-preserving. Then the set $D(a)$ defined by $\hat f(a)\leq \hat f(x)<x<a$ is nonempty (for all $a\in p(M)$). Observe that  $D(a)\subset [f(a),a]\subset p(M)$,  so  $\emptyset\neq D(b)\subset p(M)$ holds for all $b\in D(a)$. Choose $b'\in D(b)$. Then:  
$$\hat f(a)\leq \hat f(b)\leq \hat f(b')<b'<b<a$$ witnesses non-minimality of $b$ in $D(a)$. Hence no element of $D(a)$ is minimal, contradicting Lemma \ref{Lemma_basic_simple}(c). Therefore, $f$ is an automorphism of $(p(M),<)$. 

By compactness, we can modify $\hat f$ so that it is strictly increasing. In fact, in what follows we will view $\hat f$ as a definable, increasing, partial function  $\hat f\subset D\times D$. Since $\hat f$ is increasing and $\hat f(p(M))=p(M)$, $\hat f(C)\subseteq C$.


Let $\phi(x,y)$ be a formula implying $x<y$ and expressing  ``$[x,y]\subset \dom (\hat f)$ and $\hat f\restriction[x,y]$ is an order-isomorphism of $(D,<)$-intervals $[x,y]$ and $[\hat f(x),\hat f(y)]$''. Then:
$$ p(x)\cup p(y)\cup\{x<y\}\vdash \phi(x,y).$$
By compactness, there is a formula $\theta(x)\in p(x)$   such that:
\begin{equation}
\models  \forall x\forall y  (\theta(x)\land \theta(y)\land x<y\rightarrow \phi(x,y)).
\end{equation}
Since $p$ is a simple type and $\theta(x)\in p$, the set $\lnot\theta(C)$ is finite, so there is a $n_0\in \omega$ such that $\{c_i\mid i\geq n_0\}\subset \theta(C)$. For all $n\geq n_0$, by (1),  $f$ maps the (two-element) interval $[c_n,c_{n+1}]$ onto the interval $[\hat f(c_n),\hat f(c_{n+1})]$, so  
$\hat f(c_n)$ and $\hat f(c_{n+1})$ are consecutive elements of $D$, i.e $\hat f(S(c_n))=S(\hat f(c_n))$ (where $S$ is the succesor function on $(D,<)$). By induction we easily get  $\hat f(S^m(c_{n_0}))=S^m(\hat f(c_{n_0}))$ for all $m\in \omega$. Since $\hat f(c_{n_0})\in C$, there is a $k\in \mathbb Z$ with $\hat f(c_{n_0})=S^k(c_{n_0})$. Then for all $n\geq n_0$: 
$$\hat f(c_n)=\hat f(S^{n-n_0}(c_{n_0}))=    S^{n-n_0}(\hat f(c_{n_0}))=S^{n-n_0+k}(c_{n_0})=S^k(S^{n-n_0}(c_{n_0}))=S^k(c_n).$$ 
Hence  $(\hat f(x)=S^k(x))\in \tp(c_n)$ for  all $n\geq n_0$.  Since $p$ is a $C$-type, $(\hat f(x)=S^k(x))\in p(x)$.
\end{proof}
 
\begin{Lemma}\label{Lemma_x<cl(y)}
Suppose that the type $p\in S_1(T)$ is simple, as witnessed by $(C,<)$  and $(D,<)$. Let $\cl(x)=\{S^n(x)\mid n\in\mathbb Z\}$. Then   $\Pi(x,y)=p(x)\cup p(y)\cup x<\cl(y)$ has a unique completion in $S_2(T)$. 
\end{Lemma}
\begin{proof}First we {\it claim} that $\tp(a/b)$ is a $C$-type for all $(a,b)\models \Pi$. Suppose not. That is, there are $a,b\in p(M)$ with  $a<\cl(b)$, and there is a formula   $\phi(x,b)\in \tp(a/b)$ with $\phi(M,b)\cap C=\emptyset$.  Since $a\in \phi(M,b)$, we may apply Lemma \ref{Lemma_basic_simple}(c): let $a'=\min \phi(M,b)$. Then $a'\in p(M)$ and  $a'\leq a<\cl(b)$. Hence $y=\min\phi(M,x)$ relatively defines a function  $f:p(M)\to p(M)$ satisfying $f(b)=a'$. By Lemma 
\ref{Lemma_f_is_Sn},  for some integer $k$ we have $a'=S^k(b)\in\cl(b)$, contradicting $a'<\cl(b)$. Therefore, $\tp(a/b)$ is a $C$-type. 

Since $M$ is $\aleph_0$-saturated, to prove that $\Pi$ has a unique completion, it suffices to show that $\tp(a,b_1)=\tp(a,b_2)$ holds for any pair $(a,b_1),(a,b_2)$  of realizations of $\Pi$.  Without loss of generality, we will  assume $\models a<b_1\leq b_2$ and  prove that $\psi(x,b_2)\in \tp(a/b_2)$ implies $\psi(x,b_1)\in\tp(a/b_1)$. Observe that, by the above claim, both $\tp(a/b_1)$ and $\tp(a/b_2)$ are $C$-types.
Also observe that for all $m\in\mathbb Z$ the element   $S^m(b_2)$ realizes $p$, so $\models \psi(c,b_2)\leftrightarrow \psi(c,S^m(b_2))$ holds for all $c\in C$. Since $\tp(a/b_2)$ is a $C$-type  $\models \psi(a,b_2)\leftrightarrow \psi(a,S^m(b_2))$ holds, so $\psi(x,b_2)\in\tp(a/b_2)$ implies $S^m(b_2)\in\psi(a,M)$ and $\cl(b_2)\subseteq \psi(a,M)$. 
Define: 
$$D_0:=\bigcup\{[x,b_2]\mid a\leq x \mbox{ and } [x,b_2]\subseteq \psi(a,M)\}.$$ 
Clearly, $D_0$ is a convex, $ab_2$-definable subset of $\psi(a,M)$. By Lemma \ref{Lemma_basic_simple}(c), there exists $b=\min D_0$. If  $b=a$,  then we are done: $b_1\in [a,b_2]\subseteq\psi(a,M)$ implies $\psi(x,b_1)\in\tp(a/b_1)$. So suppose $a<b$. 
Then $S^{-1}(b)\notin D_0$ implies  $\models  \lnot\psi(a,S^{-1}(b))\land \psi(a,b)$. Since $\tp(S^{-1}(b))=\tp(b)$, the formula $\lnot\psi(x,S^{-1}(b))\land \psi(x,b)$ is not satisfied by an element of $C$, so $\tp(a/b)$ is not a $C$-type. 
Then  the {\it claim} implies $a\in \cl(b)$ which combined with $a<\cl(b_1)$ implies $b<b_1\leq b_2$. Since $D_0$ is convex and contains $b,b_2$, we conclude $b_1\in D_0$ and, in particular, $\psi(x,b_1)\in \tp(a/b_1)$, as desired. 
\end{proof}





\begin{Lemma}\label{Lemma_cl_n}
Suppose that $p\in S_1(T)$ is a simple type, as witnessed by $(C,<)$ and $(D,<)$.
For all finite nonempty sets $A\subset p(M)$ the partial type $\Pi_A(x)=p(x)\cup x<\cl(\min A)$ has a unique completion $p_A(x)\in S_1(A)$.  The type $p_A(x)$ is simple, as witnessed by $(C,<)$ and $(D,<)$.
\end{Lemma}
\begin{proof}First we show that the second claim follows from the first one: Note that every $C$-type from $S_1(A)$ contains the type $\Pi_A(x)$, because every formula from $\Pi_A(x)$ is satisfied by all but finitely many elements of $C$. Therefore, if $p_A(x)$ is the unique completion of $\Pi_A(x)$, then it is also the unique $C$-type in $S_1(A)$;  $p_A$ is a simple type, as witnessed by $(C,<)$.


The first claim we prove by induction on $|A|$.  Lemma \ref{Lemma_x<cl(y)} proves the case $|A|=1$. 
Suppose that the lemma holds for all simple types (in any theory) and all subsets of $p(M)$  of size at most $n$. Let $A=\{a_1,...a_{n+1}\}$ where $a_{n+1}\leq a_n\leq ...\leq a_{1}$. Then $\Pi_A(x)=p(x)\cup\{x<a_{n+1}\}$. 

Case 1. \ $a_{n+1}\in \cl(a_n)$. \ In this case $\cl(a_n)=\cl(a_{n+1})$, so the types $\Pi_{\{a_1...a_n\}}(x)=p(x)\cup x<\cl(a_n)$ and $\Pi_{A}(x)=p(x)\cup x<\cl(a_{n+1})$ 
are equal and, by the induction hypothesis, the type $\Pi_A(x)$ has a unique completion $p_{\{a_1...a_n\}}\in S_1(a_1,...,a_n)$. Since 
$a_{n+1}\in \dcl(a_n)$, the type $p_{\{a_1...a_n\}}$ has a unique extension $p_A\in S_1(A)$; $p_A$  is a unique extension of $\Pi_A(x)$, as well.   

Case 2. \ $a_{n+1}<\cl(a_n)$. \ Let $b,b'\models \Pi_A$ and we will prove $\tp(b/A)=\tp(b'/A)$. By the induction hypothesis, the type $p(x)\cup x<\cl(a_n)$ has a unique extension $p_{\{a_1,...,a_n\}}(x)\in S_1(a_1,...,a_n)$. By the first paragraph of the proof the type  $p_{\{a_1,...,a_n\}} $ is  simple, as  witnessed by $(C,<)$. Note that $b,b',a_{n+1}$ realize $\Pi_{\{a_1,...,a_n\}}(x)$, so by the induction hypothesis, they realize $p_{\{a_1.,..,a_n\}}$, too. By Lemma \ref{Lemma_x<cl(y)}  the type $p_{\{a_1,...,a_n\}}(x)\cup p_{\{a_1,...,a_n\}}(y) \cup x<\cl(y)$ has a unique completion   $q\in S_2(a_1,...,a_n)$. Note that $(b,a_{n+1}), (b',a_{n+1})$ realize $q$, so  $\tp(b/A)=\tp(b'/A)$. Therefore, the type $\tp(b/A)$ is a  unique completion of $\Pi_A(x)$, as desired.
\end{proof}
As an immediate corollary we have the following.
\begin{Corollary}\label{Cor_pntypes}
Suppose that $p\in S_1(T)$ is a simple type, as witnessed by $(C,<)$. Then for all $n\geq 1$ the type 
 \ $\Pi^n(\bar x):=\bigcup_{i=0}^n p(x_i)\cup \bigcup_{i=0}^{n-1} x_{i+1}<\cl(x_i)$ \
has a unique completion in $S_{n+1}(T)$. 
\end{Corollary}

\begin{proof}[{\it Proof of Theorem 2}]  Fix a countable, discrete, endless order $\mathbb L=(L,<_L)$ and, by saturation,  find its isomorphic copy $A_L=\{ a_i\mid i\in L\}\subset p(M)$  which is closed under the successor function of $(p(M),<)$. 
Consider the type $\Sigma(x)=p(x)\cup \{x\neq a_i\mid i\in L\}$. We {\it claim} that there is a countable model $M_{\mathbb L}$ omitting $\Sigma$. Suppose not. Then, by the omitting types theorem, there exists a consistent formula $\phi(x, \bar a)$ ($\bar a\subset A_L$) forcing $\Sigma(x)$. Choose $\phi$ with $|\bar a|=n+1$ minimal possible and, without loss of generality, let $\bar a=(a_0,...,a_n)$. 
The minimality assumption guarantees that the sets $\cl(a_i) \ ( i\leq n)$ are pairwise disjoint. Since the formula $\phi(x,\bar a)$ is consistent with $p(x)$, the set $\phi(M,\bar a)$ meets $p(M)$ and Lemma \ref{Lemma_basic_simple}(c) applies:  $b=\min \phi(M,\bar a)$ is well-defined. Since $\phi(x,\bar a)\vdash \Sigma(x)$, $b\in p(M)$ and $\cl(b)\cap \cl(a_i)=\emptyset$ for all $i\leq n$. 
Hence the sets $\cl(b), \cl(a_0),...,\cl(a_{n})$ are pairwise disjoint. Re-enumerate the tuple $a_0,...,a_n,b$ in increasing (in $p(M)$) order, and note that the new tuple satisfies the type $\Pi^{n+1}(\bar x)$ from Corollary \ref{Cor_pntypes}.  
Since $\cl(S^{-1}(b))=\cl(b)$,  the same re-enumeration applied to the tuple $a_0,...,a_n, S^{-1}(b)$ satisfies $\Pi^{n+1}(\bar x)$, too. By Corollary \ref{Cor_pntypes} we have $\tp(\bar a,b)=\tp(\bar a,S^{-1}(b))$ and, in particular, $\models  \phi(S^{-1}(b),\bar a)$. That contradicts $b=\min\phi(M,\bar a)$, proving the claim and Theorem 2 as $p(M_{\mathbb L})\cong \mathbb L$. 
\end{proof}

\begin{thebibliography}{9}
\bibitem{DI} D.Ili\'c. "Simple types in discretely ordered structures." Archive for Mathematical Logic Vol.53.7 (2014)  pp.929-947.
\end{thebibliography}
\end{document}
